\documentclass[10pt,a4paper]{amsart}
\usepackage{amsmath,amssymb,amsthm,graphicx}
% For commutative diagrams we use Xy-pic
\usepackage[all]{xy}

% We use 2cell for 2-commutative diagrams.
\xyoption{2cell}
% Theorem environments.
%
\theoremstyle{plain}
\newtheorem{theorem}[subsection]{Theorem}
\newtheorem{proposition}[subsection]{Proposition}
\newtheorem{lemma}[subsection]{Lemma}

\theoremstyle{definition}
\newtheorem{definition}[subsection]{Definition}
\newtheorem{example}[subsection]{Example}
\newtheorem{exercise}[subsection]{Exercise}
\newtheorem{situation}[subsection]{Situation}

\theoremstyle{remark}
\newtheorem{remark}[subsection]{Remark}
\newtheorem{remarks}[subsection]{Remarks}

\numberwithin{equation}{subsection}

% Macros
%
\def\lim{\mathop{\mathrm{lim}}\nolimits}
\def\colim{\mathop{\mathrm{colim}}\nolimits}
\def\Spec{\mathop{\mathrm{Spec}}}
\def\Hom{\mathop{\mathrm{Hom}}\nolimits}
\def\Ext{\mathop{\mathrm{Ext}}\nolimits}
\def\SheafHom{\mathop{\mathcal{H}\!\mathit{om}}\nolimits}
\def\SheafExt{\mathop{\mathcal{E}\!\mathit{xt}}\nolimits}
\def\Sch{\mathit{Sch}}
\def\Mor{\mathop{\mathrm{Mor}}\nolimits}
\def\Ob{\mathop{\mathrm{Ob}}\nolimits}
\def\Sh{\mathop{\mathit{Sh}}\nolimits}
\def\NL{\mathop{N\!L}\nolimits}
\def\CH{\mathop{\mathrm{CH}}\nolimits}
\def\proetale{{pro\text{-}\acute{e}tale}}
\def\etale{{\acute{e}tale}}
\def\QCoh{\mathit{QCoh}}
\def\Ker{\mathop{\mathrm{Ker}}}
\def\Im{\mathop{\mathrm{Im}}}
\def\Coker{\mathop{\mathrm{Coker}}}
\def\Coim{\mathop{\mathrm{Coim}}}

% Boxtimes
%
\DeclareMathSymbol{\boxtimes}{\mathbin}{AMSa}{"02}

%
% Macros for moduli stacks/spaces
%
\def\QCohstack{\mathcal{QC}\!\mathit{oh}}
\def\Cohstack{\mathcal{C}\!\mathit{oh}}
\def\Spacesstack{\mathcal{S}\!\mathit{paces}}
\def\Quotfunctor{\mathrm{Quot}}
\def\Hilbfunctor{\mathrm{Hilb}}
\def\Curvesstack{\mathcal{C}\!\mathit{urves}}
\def\Polarizedstack{\mathcal{P}\!\mathit{olarized}}
\def\Complexesstack{\mathcal{C}\!\mathit{omplexes}}
% \Pic is the operator that assigns to X its picard group, usage \Pic(X)
% \Picardstack_{X/B} denotes the Picard stack of X over B
% \Picardfunctor_{X/B} denotes the Picard functor of X over B
\def\Pic{\mathop{\mathrm{Pic}}\nolimits}
\def\Picardstack{\mathcal{P}\!\mathit{ic}}
\def\Picardfunctor{\mathrm{Pic}}
\def\Deformationcategory{\mathcal{D}\!\mathit{ef}}
\usepackage{cleveref}
\crefname{originalreference}{original reference}{original references}
\Crefname{originalreference}{Original reference}{Original references}
\makeatletter
\newlabel{lemma-sections-upstairs}{{0A3H}{source}{}{sourcepointer-1}{}}
\newlabel{lemma-sections-upstairs@cref}{{[originalreference][0][]0A3H}{[1][0][]source}{}{}{}}
\makeatother
\crefname{originallemma}{original lemma}{original lemmas}
\Crefname{originallemma}{Original lemma}{Original lemmas}
\makeatletter
\newlabel{more-morphisms-lemma-cover-by-geometrically-connected}{{0EY5}{source}{}{sourcepointer-2}{}}
\newlabel{more-morphisms-lemma-cover-by-geometrically-connected@cref}{{[originallemma][0][]0EY5}{[1][0][]source}{}{}{}}
\makeatother
\crefname{originallemma}{original lemma}{original lemmas}
\Crefname{originallemma}{Original lemma}{Original lemmas}
\makeatletter
\newlabel{sites-lemma-localize-morphism-strong}{{03EF}{source}{}{sourcepointer-3}{}}
\newlabel{sites-lemma-localize-morphism-strong@cref}{{[originallemma][0][]03EF}{[1][0][]source}{}{}{}}
\makeatother
\title{The pullback along a flat finitely presented morphism with geometrically reduced fibres commutes with products}
\author{The Stacks Project Authors}
\date{stacks-project:etale-cohomology.tex:local-snapshot-20261011; GNU FDL 1.2 or later}
\begin{document}
\maketitle
\noindent\textit{Editorial note.} All delivered mathematical text is the author's own, copied from the pinned Stacks Project chapter file etale-cohomology.tex; no statement or proof step was rewritten, reordered or removed.\par\medskip
\noindent\textit{Editorial note.} The counted proof is the author's complete argument: the proof is written out in full: the reduction to a good object with the required factorisation, the displayed five-term chain with the reason for each equality, and the covering argument giving the statement for all objects.\par\medskip
\noindent\textit{Editorial note.} The label lemma-fppf-reduced-fibres-pullback-products names the counted statement, and the proof environment immediately following it in the source is delivered verbatim as its proof.\par\medskip
\noindent\textit{Editorial note.} The counted statement is carried by the author's own lemma environment in the chapter source.\par\medskip
\noindent\textit{Editorial note.} The statement and proof are copied from the pinned chapter file; the three cited inputs are kept as external pointers with their canonical Stacks tags.\par\medskip
\noindent\textit{Editorial note.} The preamble is the chapter's own native preamble excerpt already used for this source: the Xy-pic setup (native preamble lines 14-18) and the native theorem declarations with the mathematical macros (native preamble lines 154-217).\par\medskip
\noindent\textit{Allowed standard prerequisites (editorial).}\begin{itemize}
\item Sections upstairs for a morphism with geometrically connected fibres. Same chapter, the lemma computing sections of a pullback on an object U whose structure morphism to the base is flat with geometrically connected fibres. Supplies the first and third equalities of the displayed chain.
\item Every etale object admits a covering by such good objects. More on Morphisms, the lemma producing an etale covering of any object by pieces that factor through an etale morphism with geometrically connected flat finitely presented fibres. Supplies the reduction from a single object to a covering, which gives the statement for all objects.
\item Locality of the pullback along a morphism of sites. Sites, the lemma identifying the restriction of f\textasciicircum{}\{-1\}G to U\_etale with the pullback along U -> V of the restriction of G. Supplies the reduction of the first and third equalities to the localised situation.
\end{itemize}\medskip
\ifdefined\hypertarget\hypertarget{sourcepointer-1}{}\fi
\noindent\textit{Source reference.} Original reference 0A3H: sources/stacks-etale-cohomology/pinned-original-source.tex, source lines 5030--5030. Prerequisite pointer invoked at proof line(s) 16771; it names a standard fact of the same chapter that the counted argument uses, and every step of the counted proof built on it is retained verbatim in the delivered body.\par\medskip
\ifdefined\hypertarget\hypertarget{sourcepointer-2}{}\fi
\noindent\textit{Source reference.} Original lemma 0EY5: sources/stacks-chow/tags.txt, source lines 18680--18680. Prerequisite pointer invoked at proof line(s) 16778; it names a standard fact of a sibling chapter that the counted argument uses, and every step of the counted proof built on it is retained verbatim in the delivered body.\par\medskip
\ifdefined\hypertarget\hypertarget{sourcepointer-3}{}\fi
\noindent\textit{Source reference.} Original lemma 03EF: sources/stacks-chow/tags.txt, source lines 4454--4454. Prerequisite pointer invoked at proof line(s) 16776; it names a standard fact of a sibling chapter that the counted argument uses, and every step of the counted proof built on it is retained verbatim in the delivered body.\par\medskip
\begin{lemma}
\label{lemma-fppf-reduced-fibres-pullback-products}
Let $f : X \to S$ be a morphism of schemes which is flat and
locally of finite presentation with geometrically reduced fibres.
Then $f^{-1} : \Sh(S_\etale) \to \Sh(X_\etale)$ commutes
with products.
\end{lemma}

\begin{proof}
Let $I$ be a set and let $\mathcal{G}_i$ be a sheaf on $S_\etale$
for $i \in I$.
Let $U \to X$ be an \'etale morphism such that $U \to S$ factors as
$U \to V \to S$ with $V \to S$ \'etale and $U \to V$ flat of finite
presentation with geometrically connected fibres. Then we have
\begin{align*}
f^{-1}(\prod \mathcal{G}_i)(U)
& =
(\prod \mathcal{G}_i)(V) \\
& =
\prod \mathcal{G}_i(V) \\
& =
\prod f^{-1}\mathcal{G}_i(U) \\
& =
(\prod f^{-1}\mathcal{G}_i)(U)
\end{align*}
where we have used Lemma \ref{lemma-sections-upstairs}
in the first and third equality
(we are also using that the restriction of $f^{-1}\mathcal{G}$
to $U_\etale$ is the pullback via $U \to V$ of the restriction of
$\mathcal{G}$ to $V_\etale$, see
Sites, Lemma \ref{sites-lemma-localize-morphism-strong}).
By More on Morphisms, Lemma
\ref{more-morphisms-lemma-cover-by-geometrically-connected}
every object $U$ of $X_\etale$ has an \'etale covering
$\{U_i \to U\}$ such that the discussion in the previous
paragraph applies to $U_i$. The lemma follows.
\end{proof}

\end{document}
